% ============================================================
%  第4章：灵敏度分析
% ============================================================
\section{第4章 \quad 灵敏度分析}

\subsection{目标系数 $c$ 的变化}

\textbf{非基变量 $c_j$ 变化：}
仅影响对应检验数 $\sigma_j = c_j - c_B^T B^{-1}a_j$。
若 $\sigma_j$ 仍 $\ge 0$（$\min$ 问题），原最优解不变。
变化范围：$c_j \le c_B^T B^{-1}a_j$。

\textbf{基变量 $c_B$ 变化：}
影响所有非基变量检验数。需重新计算 $\sigma_N = c_N - (c_B + \Delta c_B)B^{-1}N$。
若全部 $\ge 0$，原最优解不变。

\subsection{右端常数 $b$ 的变化}

$b \to b + \Delta b$ 后，基变量取值变为 $x_B = B^{-1}(b + \Delta b)$。
若 $B^{-1}(b + \Delta b) \ge 0$，则当前基仍可行且最优（目标值变为 $c_B^T B^{-1}(b+\Delta b)$）。
若出现负分量，需用对偶单纯形法恢复可行性。

\subsection{新增约束}

在原最优表基础上添加新约束行。若原最优解满足新约束，则仍为最优。
否则添加松弛/人工变量，用对偶单纯形法继续迭代。

\subsection{新增变量}

计算新变量的检验数 $\sigma_{n+1} = c_{n+1} - c_B^T B^{-1}a_{n+1}$。
若 $\sigma_{n+1} \ge 0$（$\min$ 问题），新变量不进基，原解仍最优。
若 $\sigma_{n+1} < 0$，新变量进基，继续迭代。

\subsection{本章速查}

\begin{itemize}
  \item \textbf{非基变量 $c_j$ 变}：只影响自身检验数
  \item \textbf{基变量 $c_B$ 变}：影响全部非基检验数，需重算
  \item \textbf{$b$ 变}：基变量取值变 $B^{-1}(b+\Delta b)$，不可行时用对偶单纯形
  \item \textbf{新增约束}：加入原最优表，必要时对偶单纯形
  \item \textbf{新增变量}：算检验数，负则进基
\end{itemize}
